\par\textbf{Procedimiento completo para comparar tu solución.} La condición activa de rigidez es $\delta=NL/(EA)=\delta_{max}$. Despejar área de cada material, calcular volumen y masa, y comprobar la tensión que resulta de esa área.

\begin{enumerate}

\item Área mínima elástica de A.

\[\fitmath{\mathrm{AA}=\frac{12000\cdot 1.2}{7\times10^{10}\cdot 0.0006}\approx 0.0003428571\,\mathrm{m^{2}}}\]

\item Área de B con su propio módulo.

\[\fitmath{\mathrm{AB}=\frac{12000\cdot 1.2}{1.2\times10^{11}\cdot 0.0006}\approx 0.0002\,\mathrm{m^{2}}}\]

\item Densidad por volumen de A.

\[\fitmath{\mathrm{mA}=2700\cdot \mathrm{AA}\cdot 1.2}\]

\[\fitmath{\mathrm{mA}=2700\cdot (0.0003428571429)\cdot 1.2\approx 1.110857\,\mathrm{kg}}\]

\item Densidad por volumen de B.

\[\fitmath{\mathrm{mB}=1600\cdot \mathrm{AB}\cdot 1.2}\]

\[\fitmath{\mathrm{mB}=1600\cdot (0.0002)\cdot 1.2\approx 0.384\,\mathrm{kg}}\]

\item Tensión resultante de A.

\[\fitmath{\mathrm{sA}=\frac{12000}{\mathrm{AA}}}\]

\[\fitmath{\mathrm{sA}=\frac{12000}{(0.0003428571429)}\approx 3.5\times10^{7}\,\mathrm{Pa}}\]

\item Tensión resultante de B.

\[\fitmath{\mathrm{sB}=\frac{12000}{\mathrm{AB}}}\]

\[\fitmath{\mathrm{sB}=\frac{12000}{(0.0002)}\approx 6\times10^{7}\,\mathrm{Pa}}\]

\item Reducción de masa bajo el único criterio de rigidez.

\[\fitmath{\mathrm{saving}=1-\frac{\mathrm{mB}}{\mathrm{mA}}}\]

\[\fitmath{\mathrm{saving}=1-\frac{(0.384)}{(1.110857143)}\approx 0.654321\,}\]

\end{enumerate}
